%==================================================================
% BAB 1 PENDAHULUAN
% Struktur mengikuti Pedoman 6.1 (Kerangka Proposal Tugas Akhir)
%  1.1 Latar Belakang
%  1.2 Perumusan Masalah
%  1.3 Tujuan
%  1.4 Manfaat Tugas Akhir
%  1.5 Batasan Masalah
%==================================================================
\chapter{PENDAHULUAN}
\label{chap:pendahuluan}

%------------------------------------------------------------------
\section{Latar Belakang}
%------------------------------------------------------------------

Indonesia merupakan negara penghasil ikan yang besar, sehingga mutu
produk perikanan menjadi faktor penting dalam rantai perdagangan.
Ikan termasuk bahan pangan yang sangat mudah mengalami kemunduran mutu
setelah ditangkap. Menurut \citet{estiasih2015}, proses kemunduran mutu
ikan berlangsung melalui tahap \textit{rigor mortis}, penguraian
senyawa oleh enzim, dan penguraian oleh aktivitas mikroorganisme.
Ketiga tahap tersebut berjalan cepat, terutama bila ikan tidak segera
didinginkan.

Penilaian kesegaran ikan di pasar tradisional umumnya dilakukan secara
organoleptik, yaitu dengan mengamati warna insang, mata, tekstur, dan
bau. Cara ini membutuhkan tenaga terlatih, bersifat subjektif, dan sulit
dilakukan dalam jumlah besar. Beberapa penelitian menunjukkan bahwa
penilaian berbasis citra digital dapat menjadi alternatif yang lebih
konsisten \citep{prasetyo2021}. Citra mata dan insang ikan menyimpan
informasi warna yang berkorelasi dengan tingkat kesegaran, sehingga
dapat diolah secara kuantitatif \citep{gonzalez2018}.

Pengolahan citra digital telah banyak dimanfaatkan untuk pemeriksaan
mutu bahan pangan secara tidak merusak. Namun, pemanfaatan teknologi
tersebut di lingkungan pasar ikan Banyuwangi belum banyak dilakukan,
padahal kebutuhan akan alat bantu penilaian yang cepat dan murah cukup
tinggi. Berdasarkan uraian tersebut, penelitian ini merancang bangun
sistem deteksi kesegaran ikan berbasis pengolahan citra digital dan
jaringan saraf tiruan yang diharapkan dapat membantu pembeli maupun
pedagang dalam menentukan mutu ikan secara cepat dan objektif.

%------------------------------------------------------------------
\section{Perumusan Masalah}
%------------------------------------------------------------------

Berdasarkan latar belakang di atas, rumusan masalah pada penelitian ini
adalah sebagai berikut.
\begin{packed_enum}
  \item Bagaimana merancang sistem akuisisi citra mata dan insang ikan
        yang menghasilkan gambar dengan pencahayaan konsisten?
  \item Bagaimana mengekstraksi ciri warna citra ikan pada ruang warna
        \textit{Hue, Saturation, Value} (HSV) sehingga dapat membedakan
        tingkat kesegaran ikan?
  \item Bagaimana kinerja jaringan saraf tiruan dalam mengklasifikasikan
        tingkat kesegaran ikan berdasarkan ciri warna yang diperoleh?
\end{packed_enum}

%------------------------------------------------------------------
\section{Tujuan}
%------------------------------------------------------------------

Tujuan yang ingin dicapai pada penelitian ini adalah:
\begin{packed_enum}
  \item merancang bangun sistem deteksi kesegaran ikan berbasis
        pengolahan citra digital;
  \item mengimplementasikan jaringan saraf tiruan untuk
        mengklasifikasikan tingkat kesegaran ikan; dan
  \item mengukur tingkat akurasi sistem yang dirancang pada data uji
        yang belum pernah digunakan pada proses pelatihan.
\end{packed_enum}

%------------------------------------------------------------------
\section{Manfaat Tugas Akhir}
%------------------------------------------------------------------

Manfaat yang diharapkan dari penelitian ini adalah sebagai berikut.
\begin{packed_enum}
  \item Bagi pembeli dan pedagang ikan, tersedia alat bantu penilaian
        kesegaran yang cepat, murah, dan tidak merusak produk.
  \item Bagi Program Studi \prodi, menjadi contoh penerapan pengolahan
        citra digital dan pembelajaran mesin pada permasalahan nyata.
  \item Bagi peneliti selanjutnya, menjadi rujukan awal pengembangan
        sistem sejenis dengan cakupan jenis ikan yang lebih luas.
\end{packed_enum}

%------------------------------------------------------------------
\section{Batasan Masalah}
%------------------------------------------------------------------

Agar pembahasan terarah, penelitian ini dibatasi pada hal-hal berikut.
\begin{packed_enum}
  \item Objek penelitian adalah ikan bandeng (\textit{Chanos chanos})
        yang diperoleh dari Pasar Banyuwangi.
  \item Ciri yang digunakan terbatas pada ciri warna pada ruang HSV,
        tidak mencakup ciri bentuk maupun tekstur.
  \item Pengambilan citra dilakukan di dalam kotak akusisi dengan
        pencahayaan buatan yang tetap.
  \item Klasifikasi dibagi menjadi tiga kelas, yaitu segar, cukup segar,
        dan tidak segar, mengacu pada kriteria organoleptik.
\end{packed_enum}
